\begin{summarybox}
\paragraph{核心：极值由附近大小决定，导数为零只是候选条件}
\[
\text{内点极值且可导}\Rightarrow f'(x_0)=0.
\]
在连续及两侧可导的前提下：
\[
+\to-:\text{极大值},\qquad -\to+:\text{极小值}.
\]
先确定定义域，再列驻点和不可导点，最后检查附近大小。
核对：点与值是否区分？是否遗漏尖点？是否误把端点最值当内部极值？变号判定的连续性前提是否成立？
\end{summarybox}
